\section*{Conclusion}

L'analyse approfondie de la relation entre les empereurs et la foi chrétienne, et la vision qu'en a Ambroise, nous a permis de mieux comprendre ce qu'il conçoit comme étant la monarchie chrétienne idéale. L'objectif principal de ce chapitre était de se détacher de la vision utopique pour confronter sa théologie politique à la réalité du pouvoir que l'évêque côtoie. Il s'agissait d'observer comment il aborde l’exercice de l'autorité et les actions impériales lorsque celles-ci touchent à la religion. Nous avons ainsi pu mettre en évidence deux aspects importants : d'une part, la compréhension des méthodes de réflexion d’Ambroise à travers les Écritures, la philosophie et la morale chrétienne ; d'autre part, l'analyse, à partir de cette approche, de la mesure de l'influence d'Ambroise sur le pouvoir romain.

\bigskip

Il apparaît assez clairement que l’évêque de Milan tire profit du contexte politico-religieux de son époque pour interférer dans les affaires de la cour, bien qu’il peine à imposer purement et simplement son idéologie sur les empereurs. Des souverains comme Gratien et Théodose se montrent particulièrement entreprenants dans la lutte contre le polythéisme, sans que l'on ait nécessairement connaissance d'une demande d'Ambroise ou d'une influence par ses travaux. Le \latin{De fide}, adressé à Gratien, se concentre par exemple davantage sur le débat autour de la doctrine de Nicée, et les premières correspondances avec Théodose sont plus tardives. L’implication religieuse de ces empereurs relève donc, en grande partie, de leur propre initiative. Mais cette dynamique lui permet de s'appuyer sur une base solide de défense de la foi pour mettre en avant sa théorie de la Providence divine et l’obligation, pour un pouvoir chrétien, de se soumettre aux desseins de Dieu. Sans un pouvoir ancré dans la foi, il est très probable que l'évêque, siégeant dans la capitale impériale, ne se serait pas permis une telle approche.

\bigskip

Soutenu par ce cadre politique favorable, la pensée politique ambrosienne s'en retrouve plus riche et audacieuse. L'ensemble de ce chapitre s'est concentré sur les exigences des vertus chrétiennes mises en avant par Ambroise, qu'elles soient formulées dans des traités théoriques, ou directement adressées aux souverains. Pour Ambroise, le véritable enjeu reste la concorde de la société qui mène à la réussite de chacun, un idéal qui ne peut être atteint qu'à travers le règne d'un monarque encadré par la morale chrétienne. Ce lien au christianisme est exprimé par Ambroise dans deux axes. Premièrement, la foi personnelle, qui permet d'inscrire dans l'inconscient des empereurs les valeurs de clémence et de modération pour établir une plus grande justice dans les actions. Deuxièmement, la lutte active pour la protection de la foi chrétienne, qui se traduit par exemple par la législation contre le polythéisme ainsi que par la force déployée contre les usurpateurs.

\bigskip

En définitive, Ambroise parvient tout de même, par sa théorie politique, à modeler en partie les empereurs romains selon ses propres volontés, s'appuyant sur des évidences théologiques pour influencer les actions impériales. Cette volonté d'encadrer le pouvoir et donc d'imposer sa vision sur le monde politique romain nous conduit alors vers un troisième chapitre, non plus concentré sur l'influence théorique et religieuse de l'évêque, mais sur son autorité politique. Nous chercherons à expliquer l’ensemble des actions concrètes qu'il met en place pour confronter ou contester le pouvoir établi.
